\chapter{描述统计}\label{ch:g11:stat}

统计将一列数压缩为几个有意义的概括量：
一是数据的集中位置（平均数、中位数），二是数据的分散程度（四分位数、
方差、标准差）。本章以同一个例子贯穿全部计算。
同样的概括量在随机变量中再次出现，
见\cref{ch:g11:prob}，其概率规律见\cref{ch:g12:sums}。

本章通篇使用的例子数据，是一个班级（满分 $20$ 分）的成绩，
该班共 $12$ 名学生：
\[
8,\ 9,\ 10,\ 10,\ 11,\ 12,\ 12,\ 12,\ 13,\ 14,\ 16,\ 17 .
\]

\section{平均数与中位数}

\begin{definition}[平均数]\label{def:g11:stat:mean}
数值 $x_1, x_2, \dots, x_n$ 的\emph{平均数}\index{平均数}为
\[
\bar x = \frac{x_1 + x_2 + \dots + x_n}{n}.
\]
若数值 $x_i$ 的频率为 $n_i$（且
$n_1 + \dots + n_k = n$），则
$\bar x = \frac{n_1 x_1 + \dots + n_k x_k}{n}$。
\end{definition}

\begin{definition}[中位数]\label{def:g11:stat:median}
对于已排序的数据集，\emph{中位数}\index{中位数}是将数据平分为两半的一个值：
即至少有一半的数据 $\leq$ 该值，且至少有一半的数据 $\geq$ 该值。
实际计算时，先将 $n$ 个数据排序；若 $n$ 为奇数，则中位数为
中间的值；若 $n$ 为偶数，则取中间两个值的
中点。
\end{definition}

\begin{example}\label{ex:g11:stat:meanmedian}
对于上述数据，总和为 $144$，故 $\bar x = \frac{144}{12} =
12$。排序后，中间两个值（第 $6$ 个与第 $7$ 个）均为
$12$，故中位数同样为 $12$。平均数与中位数不必相等：
中位数不考虑极端值\emph{离中心有多远}，而平均数则考虑。
将最高分 $17$ 换成 $20$，平均数变为 $12.25$，而
中位数仍为 $12$。
\end{example}

\section{四分位数}

\begin{definition}[四分位数、四分位距]\label{def:g11:stat:quartiles}
\emph{第一四分位数}\index{四分位数} $Q_1$ 是满足
至少四分之一的数据 $\leq Q_1$ 的最小值；\emph{第三
四分位数} $Q_3$ 是满足至少四分之三的数据 $\leq Q_3$ 的最小值。
\emph{四分位距}\index{四分位距}为 $Q_3 - Q_1$，即
数据中间一半的宽度。
\end{definition}

\begin{example}
$n = 12$：四分之一为 $3$，故 $Q_1$ 为排序后的第 $3$ 个值：
$Q_1 = 10$。四分之三为 $9$，故 $Q_3$ 为第 $9$ 个值：
$Q_3 = 13$。四分位距为 $3$：该班中间一半的学生
成绩集中在 $3$ 分的范围内。
\end{example}

\begin{omfigure}
\begin{tikzpicture}[x=0.62cm]
  \draw[->] (7,0) -- (18.4,0) node[right, font=\small] {分数};
  \foreach \x in {8,10,12,14,16} {
    \draw (\x,0) -- (\x,-0.12) node[below, font=\small] {\x};
  }
  \fill[omDef] (8,0.25) circle (2pt);
  \fill[omDef] (9,0.25) circle (2pt);
  \fill[omDef] (10,0.25) circle (2pt);
  \fill[omDef] (10,0.55) circle (2pt);
  \fill[omDef] (11,0.25) circle (2pt);
  \fill[omDef] (12,0.25) circle (2pt);
  \fill[omDef] (12,0.55) circle (2pt);
  \fill[omDef] (12,0.85) circle (2pt);
  \fill[omDef] (13,0.25) circle (2pt);
  \fill[omDef] (14,0.25) circle (2pt);
  \fill[omDef] (16,0.25) circle (2pt);
  \fill[omDef] (17,0.25) circle (2pt);
  \draw[omThm, thick] (12,-0.35) -- (12,1.15)
    node[above, font=\small] {$\bar x = 12$};
  \draw[omProp, thick] (10,1.55) -- (10,2.35) -- (13,2.35) -- (13,1.55)
    -- cycle;
  \draw[omProp, thick] (12,1.55) -- (12,2.35);
  \draw[omProp, thick] (8,1.95) -- (10,1.95);
  \draw[omProp, thick] (13,1.95) -- (17,1.95);
  \draw[omProp, thick] (8,1.75) -- (8,2.15);
  \draw[omProp, thick] (17,1.75) -- (17,2.15);
  \node[omProp, font=\small, above] at (10,2.35) {$Q_1$};
  \node[omProp, font=\small, above] at (13,2.35) {$Q_3$};
\end{tikzpicture}

{\small 本例数据的点图，其平均数（红色）以及
其上方的五数概括（橙色）：最小值、$Q_1$、中位数、$Q_3$、
最大值。方框包含数据中间一半的部分。}
\end{omfigure}

\section{方差与标准差}

\begin{definition}[方差、标准差]\label{def:g11:stat:variance}
数据集的\emph{方差}\index{方差}是
各数据与平均数之差平方的平均数：
\[
V = \frac{1}{n}\sum_{i=1}^{n} (x_i - \bar x)^2,
\]
而\emph{标准差}\index{标准差}为
$\sigma = \sqrt{V}$，并与数据具有相同的单位。
\end{definition}

\begin{proposition}[简便公式]\label{prop:g11:stat:shortcut}
方差等于各数值平方的平均数减去平均数的平方：
\[
V = \frac{1}{n}\sum_{i=1}^{n} x_i^2 \;-\; \bar x^{\,2}.
\]
\end{proposition}

\begin{proof}
展开每个平方偏差：
$(x_i - \bar x)^2 = x_i^2 - 2\bar x\, x_i + \bar x^2$。求和并
除以 $n$：
\[
V = \frac1n \sum x_i^2 - 2\bar x \cdot \underbrace{\frac1n \sum
x_i}_{=\ \bar x} + \frac1n \cdot n\bar x^2
= \frac1n\sum x_i^2 - 2\bar x^2 + \bar x^2
= \frac1n\sum x_i^2 - \bar x^2 .
\]
对随机变量而言，同一恒等式为
\cref{prop:g11:prob:konig}。
\end{proof}

\begin{example}\label{ex:g11:stat:variance}
对于本例数据，与 $\bar x = 12$ 的偏差依次为 $-4$，
$-3$、$-2$、$-2$、$-1$、$0$、$0$、$0$、$1$、$2$、$4$、$5$；其平方
和为
\[
16 + 9 + 4 + 4 + 1 + 0 + 0 + 0 + 1 + 4 + 16 + 25 = 80 .
\]
由此
\[
V = \frac{80}{12} = \frac{20}{3} \approx 6.67,
\qquad
\sigma = \sqrt{20/3} \approx 2.6 \text{ 分}.
\]
粗略地说，一个典型分数距
平均数约 $2.6$ 分。
\end{example}

\begin{proposition}[单位变换]\label{prop:g11:stat:affine}
若每个数值经变换 $y_i = a x_i + b$，则
\[
\bar y = a \bar x + b,
\qquad
V_y = a^2\, V_x,
\qquad
\sigma_y = \abs{a}\,\sigma_x .
\]
\end{proposition}

\begin{proof}
$\bar y = \frac1n \sum (a x_i + b) = a\left(\frac1n\sum x_i\right) +
\frac1n \cdot nb = a\bar x + b$。于是
$y_i - \bar y = a(x_i - \bar x)$：每个偏差都乘以 $a$
（平移量 $b$ 相消），因此每个平方偏差都乘以 $a^2$，
其平均数亦然。开方得
$\sigma_y = \abs a\,\sigma_x$。
\end{proof}

\begin{example}
以摄氏度记录的温度，平均数为 $20$，标准
差为 $3$，经 $F = 1.8\,C + 32$ 换算为华氏度：平均数
变为 $1.8 \times 20 + 32 = 68$，标准差
变为 $1.8 \times 3 = 5.4$。平移数据集不改变其离散程度；
对数据作伸缩，则离散程度随之伸缩。
\end{example}

\begin{method}[概括与比较数据集]\label{met:g11:stat:compare}
比较两个数据集时，可并列列出：平均数（或中位数）用于
\emph{位置}，标准差（或四分位距）用于
\emph{离散程度}。（中位数，四分位距）这一对量对
极端值具有稳健性；（平均数，标准差）这一对量则利用了每一个数值，
并为概率论提供基础。
\end{method}

\begin{example}
两名射手各射 $10$ 支箭；两人的得分平均数均为 $8.5$，但
标准差分别为 $0.5$ 和 $2.1$。平均水平相同 --- 但
第一名射手的发挥明显更稳定。
\end{example}

\section{练习}

\begin{exercise}[$\star$]\label{exo:g11:stat:1}
求下列数据集的平均数、中位数、四分位数及
四分位距：
\[
5,\ 7,\ 7,\ 8,\ 9,\ 10,\ 11,\ 14 .
\]
\end{exercise}

\begin{exercise}[$\star$]\label{exo:g11:stat:2}
求下列数据集的方差与标准差
$2,\ 4,\ 4,\ 4,\ 5,\ 5,\ 7,\ 9$（先求平均数，再用
\cref{prop:g11:stat:shortcut}给出的简便公式）。
\end{exercise}

\begin{exercise}[$\star$]\label{exo:g11:stat:3}
掷一颗骰子 $50$ 次，结果用频率汇总如下：
\begin{center}
\begin{tabular}{c|cccccc}
结果 & 1 & 2 & 3 & 4 & 5 & 6\\
\hline
频率 & 6 & 9 & 8 & 10 & 9 & 8
\end{tabular}
\end{center}
求结果的平均数与中位数。
\end{exercise}

\begin{exercise}[$\star\star$]\label{exo:g11:stat:4}
一个班级有 $15$ 名学生，平均分为 $11.2$。又有第 $16$ 名学生加入，
得分为 $18$。求新班级的平均分。
\end{exercise}

\begin{exercise}[$\star\star$]\label{exo:g11:stat:5}
数据集 $6, 8, x, 12, 14$ 的平均数为 $10$。求 $x$，再求
补全后数据集的标准差。
\end{exercise}

\begin{exercise}[$\star$]\label{exo:g11:stat:6}
一组成绩的平均数为 $12.4$，标准差为 $2.5$，现经变换
$y = \frac{x}{2} + 4$。求新的平均数与标准差。
\end{exercise}

\begin{exercise}[$\star\star$]\label{exo:g11:stat:7}
A 组（$20$ 个数值）的平均数为 $10$；B 组（$30$ 个数值）的平均数为 $15$。
求合并后全部 $50$ 个数值的平均数，它是否为中点
$12.5$？若不是，请说明理由。
\end{exercise}

\begin{exercise}[$\star\star$]\label{exo:g11:stat:8}
两条生产线灌装 $1$ 千克装的糖袋。抽样结果如下：
生产线 1：平均数 $1.002$ 千克，标准差 $0.006$ 千克；生产线 2：平均数
$1.001$ 千克，标准差 $0.019$ 千克。应先重新校准哪一条生产线，
并说明理由。
\end{exercise}

\begin{exercise}[$\star\star$]\label{exo:g11:stat:9}
一组含 $n = 10$ 个数值的数据满足 $\sum x_i = 70$ 且
$\sum x_i^2 = 540$。求其平均数、方差与标准差。
\end{exercise}

\begin{exercise}[$\star\star\star$]\label{exo:g11:stat:10}
数据集 $1,\ 2,\ 2,\ 3,\ 3,\ 3,\ 4,\ 4,\ 5,\ 23$ 中含有一个
\emph{离群值}。
\begin{enumerate}
  \item 分别计算包含与不包含 $23$ 时的平均数与中位数。
  \item 分别计算包含与不包含 $23$ 时的四分位距。
  \item 哪些统计量对离群值具有稳健性？
        并据此给出建议。
\end{enumerate}
\end{exercise}

\section{问题：标准分数，一把通用的尺子}

\begin{problem}[{周末问题 --- 标准差比较苹果与橙子，
切比雪夫不等式把 $\sigma$ 变成一种保证，
而平均数赢得自己的奖牌}]
\label{pb:g11:stat:1}
在难卷上得 $15/20$ 是否比在易卷上得 $13/20$ 更好？
一列较小的数中出现的 $23$ 是偶然还是离群值？原始
数值无法回答这些问题——但除以恰当的标准差，
它们便说同一种语言，即\emph{标准分数}。本题
练习方差的计算方法，
（\cref{def:g11:stat:variance}、
\cref{prop:g11:stat:shortcut}），构造通用的尺子，并
证明一个著名的不等式，它使 $\sigma$ 由描述性的数字
变为一种保证。

\medskip
\noindent\textbf{第一部分 --- 方差的计算。}
\begin{enumerate}
  \item 对数据集 $4, 8, 6, 5, 2$，先求平均数，再
        由定义直接求方差，然后
        求 $\sigma$。
  \item 用简便公式重新计算方差
        （\cref{prop:g11:stat:shortcut}）——平方的平均数
        减去平均数的平方，并检验两者结果一致。
  \item 已知频数：成绩为 $8$、$10$、$14$ 的学生分别有
        $3$、$5$、$2$ 名，求平均数、方差与
        $\sigma$（保留两位小数）。
  \item 气温的平均数为 $20$，$\sigma = 3$（摄氏度）。求以
        华氏度表示的平均数与 $\sigma$
        （$F = 1.8\,C + 32$；\cref{prop:g11:stat:affine}）。
  \item 证明刚才用到的两条规则：给每个数值加上常数 $b$，
        平均数增加 $b$，$\sigma$ 保持不变；
        给每个数值乘以 $a$，平均数乘以 $a$，
        $\sigma$ 乘以 $\abs a$。
\end{enumerate}

\medskip
\noindent\textbf{第二部分 --- 通用的尺子。}
数据集中某数值的\emph{标准分数}\index{标准分数}，记该数值为 $x$，
数据集平均数为 $\bar x$、标准差为 $\sigma$，则
$z = \dfrac{x - \bar x}{\sigma}$，它表示 $x$ 与平均数
相差的标准差个数。
\begin{enumerate}[resume]
  \item 考试 A 的平均数为 $12$，$\sigma = 2$；考试 B 的平均数为
        $8$，$\sigma = 4$。爱丽丝在 A 中得 $15$ 分，鲍勃
        在 B 中得 $13$ 分。分别计算两人的标准分数，谁的成绩
        高出平均水平更多？
  \item 求第 1 题中数据集的五个标准分数，
        并验证它们的平均数为 $0$、标准差为
        $1$。由第 5 题的结论说明
        为何标准化\emph{总是}得到平均数为 $0$、
        $\sigma = 1$。
  \item 某能力测验的平均数为 $100$，$\sigma = 15$。原始
        分数为多少时 $z = 2$？$76$ 分的标准分数是多少？
  \item 儿科医生记录某婴儿的体重为
        “$z = -2.5$”。请说明其含义，以及为何
        生长状况要用标准分数而非千克来追踪，
        尤其是随孩子年龄增长。
  \item \cref{exo:g11:stat:10}中的离群值（数据集
        $1, 2, 2, 3, 3, 3, 4, 4, 5, 23$）：计算平均数、
        $\sigma$ 以及 $23$ 的标准分数。经典的警戒阈值
        为 $\abs z \geq 3$，结论如何？
  \item 重新计算 $\sigma$（\emph{不含} $23$ 的数据），
        并求 $23$ 按干净数据的平均数与
        $\sigma$ 算得的标准分数。离群值对本用来度量它的
        尺子产生了什么影响？
        由\cref{exo:g11:stat:10}可知应采取什么对策？
\end{enumerate}

\medskip
\noindent\textbf{第三部分 --- 切比雪夫的保证。}
\begin{enumerate}[resume]
  \item 证明数据集的比内梅–切比雪夫不等式：
        若比例为 $p$ 的数值满足
        $\abs{x - \bar x} \geq k\sigma$，则在方差和式中
        仅保留这些项，可得
        \[
        \sigma^2 \geq p \cdot (k\sigma)^2,
        \qquad\text{于是}\qquad
        p \leq \frac{1}{k^2} .
        \]
  \item 该不等式对任一数据集中偏离平均数超过 $2\sigma$ 的比例
        能保证什么？超过 $3\sigma$ 的比例又如何？
        请在第 1 题的数据集上检验关于
        $2\sigma$ 的结论。
  \item 某工厂生产螺栓，平均数为 $50$ 毫米，
        $\sigma = 0.1$ 毫米。切比雪夫不等式对落在
        $\intcc{49.7}{50.3}$ 之外的螺栓比例能给出什么保证？
        （实际生产远好于此，
        ——十二年级学的钟形曲线会把这个保守的
        $\frac19$ 化为不到百分之一的比例。）
  \item 两只基金的年平均收益率均为 $5\,\%$，
        其中基金 A 的 $\sigma = 2\,\%$，基金 B 的 $\sigma = 12\,\%$。
        对每只基金，切比雪夫不等式对
        收益率低于某一水平的年份所占比例（可视为概率）有何说明？
        该水平为 $-19\,\%$。$\sigma$ 在金融中通常叫作什么？
\end{enumerate}

\medskip
\noindent\textbf{第四部分 --- 平均数的奖章与一则警告。}
\begin{enumerate}[resume]
  \item 设 $f(c)$ 为离差平方
        $(x_i - c)^2$ 的平均数，中心 $c$ 任取。证明
        恒等式
        \[
        f(c) = \sigma^2 + (\bar x - c)^2 ,
        \]
        并由此断言：平均数是唯一使平均离差平方
        最小的点 --- 它的奖章，与中位数使绝对离差
        最小的性质相呼应，见
        \cref{pb:g10:stats:1}。
  \item 以第 1 题的数据验证恒等式，
        取 $c = 6$，分别直接计算 $f(6)$ 与用公式计算。
  \item 同样的统计量、不同的数据：验证
        $P = \{3, 5, 7, 9, 11\}$ 与 $Q = \{1, 7, 7, 7, 13\}$
        的平均数相同而方差不同；然后
        构造形如
        $\{7 - a, 7, 7, 7, 7 + a\}$ 的数据集，
        使其平均数与方差均与 $P$ 恰好相同，却自有一种形状。
  \item 证明：若 $\sigma = 0$，则所有值都相等。（
        零平方和的每一项须为何值？）
  \item 终章 --- 统计学家的阶梯，每一级一行：
        中心（平均数、中位数），离散程度（$\sigma$、四分位距），
        通用标尺（$z$），保证（切比雪夫不等式），
        以及第 18 题的警告：统计量会压缩数据，
        而压缩必然丢失信息 --- 十二年级的钟形曲线
        正等待把标准分数化为精确的概率。
\end{enumerate}
\end{problem}
